% Generated draft: requires mathematical and source review.
\documentclass[11pt]{article}
\usepackage[margin=25mm]{geometry}
\usepackage{fontspec}
\setmainfont{Latin Modern Roman}
\usepackage{amsmath,amssymb,mathtools,unicode-math}
\setmathfont{Latin Modern Math}
\usepackage{enumitem,needspace,xurl,hyperref}
\setlength{\parindent}{0pt}
\setlength{\parskip}{5pt}
\setlength{\emergencystretch}{3em}
\newcommand{\N}{\mathbb{N}}
\newcommand{\Z}{\mathbb{Z}}
\newcommand{\Q}{\mathbb{Q}}
\newcommand{\R}{\mathbb{R}}
\newcommand{\C}{\mathbb{C}}
\newcommand{\F}{\mathbb{F}}
\newcommand{\A}{\mathbb{A}}
\newcommand{\PP}{\mathbb{P}}
\newcommand{\E}{\mathbb{E}}
\newcommand{\eps}{\varepsilon}
\newcommand{\dd}{\,\mathrm{d}}
\begin{document}
\section*{UnsolvedMath ID: 20000812}
\textbf{Source problem code:} catalogue entry
\subsection{Definitions and mathematical statement}
The source asks to describe the irreducible component(s) of \(E_\rho\), called the moduli of finite-length graded modules with function \(\rho\). Interpreting \(\rho\) as a Hilbert function means \(\rho(d)=\dim_kM_d\), with finite support. The original sentence does not specify the ground field, graded ring, scheme or stack construction of \(E_\rho\), or a distinguished component; its singular word ``component'' does not by itself select one.

The supplied partial result uses an algebraically closed field \(k\), the standard graded ring \(S=k[x_1,x_2,x_3,x_4]\), and \(\rho=(p,1,r)\) in degrees \(0,1,2\), zero elsewhere, with \(p,r\ge1\). Fix bases of
\[
 M_0=k^p,\qquad M_1=k,\qquad M_2=k^r.
\]
Write multiplication by \(x_i\) from \(M_0\) to \(M_1\) as a row \(a_i\), and from \(M_1\) to \(M_2\) as a column \(b_i\). Let \(X\) be the \(4\)-by-\(p\) matrix with rows \(a_i\), and \(Y\) the \(4\)-by-\(r\) matrix with rows \(b_i^{\mathsf T}\). The framed module scheme is defined by the cross-commutation ideal
\[
 J=(b_i a_j-b_j a_i:1\le i<j\le4),
\]
where each matrix equation contributes all its entries to the polynomial ideal in the entries of \(X,Y\).

The aid reports the exact radical decomposition
\[
 \sqrt J=(X)\cap(Y)\cap I_2([X\ Y]),
\]
where \((X)\), \((Y)\) are the ideals of all entries of the respective matrices and \(I_2\) is the ideal of \(2\)-by-\(2\) minors. Equivalently the reduced framed scheme is the union of
\[
 A=\{X=0\},\qquad B=\{Y=0\},\qquad
 C=\{\operatorname{rank}[X\ Y]\le1\}.
\]
Their dimensions are \(\dim A=4r\), \(\dim B=4p\), and \(\dim C=p+r+3\). Removing contained loci, the reported irreducible components are: only \(C\) if \(p=r=1\); \(A,C\) if \(p=1<r\); \(B,C\) if \(r=1<p\); and \(A,B,C\) if \(p,r>1\).

At a generic mixed point of \(C\), both adjacent multiplications are nonzero and have a common variable direction: after linear changes, the \(x_i\) act as scalar multiples of one degree-one operator. The corresponding generic graded module is a length-three one-variable chain with one vector in each of degrees \(0,1,2\), together with \(p-1\) annihilated vectors in degree \(0\) and \(r-1\) annihilated vectors in degree \(2\). On \(A\) the first adjacent multiplication vanishes, and on \(B\) the second vanishes. The supplied summary says generic types are proved in accompanying artifacts, but those artifacts are not included here. These conclusions concern the reduced framed scheme; they do not supply the full scheme structure, classify arbitrary \(\rho\), or settle what moduli interpretation of \(E_\rho\) the original problem intends.
\subsection{Short English statement}
The broad question asks for components of a moduli space of finite-length graded modules. The supplied aid treats the specific three-degree Hilbert function \((p,1,r)\) over four variables and reports a reduced framed-space decomposition into one, two, or three components.
\subsection{Sources}
\begin{itemize}
\item Local UnsolvedMath registry, permanent ID \(20000812\), source code \texttt{catalogue entry}; \texttt{aim-arithmetic-geometry-notes.json}, \texttt{notes[57]}, tag \texttt{problem}.
\item AIM workshop ``Components of Hilbert schemes,'' Problem 11: \url{https://aimath.org/WWN/hilbertschemes/hilbertschemes.pdf}. Statement status \texttt{exact}. No original remarks, literature references, proposer, or year are supplied. No PDF was accessed for this draft.
\item Supplied research attempt 1: \texttt{partial\_result}, problem status at run \texttt{open}, rigor label \texttt{proved}, novelty confidence low; registry status \texttt{partially\_solved}, classification \texttt{PARTIAL-PROGRESS}. Created and updated 14 August 2026. No authoritative external citation for the proposed radical decomposition is supplied. The machine-generated aid is not peer reviewed and requires independent verification before citation.
\item Ulam AI / UnsolvedMath Contributors; CC BY 4.0.
\end{itemize}
\end{document}
